\begin{definition}[The two-plaquette block at an arbitrary torus position]%
    \label{def:peps_translated_two_plaquette_geometry}
    \lean{TNLean.PEPS.translatedTwoPlaquetteVertex,
        TNLean.PEPS.translatedTwoPlaquetteRegion,
        TNLean.PEPS.translatedTwoPlaquetteVertexEquiv,
        TNLean.PEPS.translatedTwoPlaquetteIso,
        TNLean.PEPS.translatedTwoPlaquetteTree,
        TNLean.PEPS.translatedTwoPlaquetteCycleBond}
    \leanok
    \uses{def:peps_two_plaquette_graph,def:peps_two_plaquette_torus_geometry,
        thm:peps_edge_map_subgraph}
    On a finite torus of horizontal period at least four and vertical period
    at least three, fix any position $v$. Translate the six-site numbering
    of the two-plaquette block by $v$, with both coordinates taken modulo
    their periods. Let $R_v$ be the resulting set of sites. Its induced graph
    is identified with the fixed two-plaquette graph. Transport the fixed
    spanning path and its two remaining bonds through that identification.
    These constructions permit either periodic seam to pass through the block.
    They describe the local geometry in
    \cite[Theorem~6.16]{Schuch2010PEPS}.
\end{definition}

\begin{theorem}[Actual translated two-plaquette geometry]%
    \label{thm:peps_translated_two_plaquette_geometry}
    \lean{TNLean.PEPS.translatedTwoPlaquetteTree_isTree,
        TNLean.PEPS.translatedTwoPlaquetteTree_le,
        TNLean.PEPS.translatedTwoPlaquetteRegion_card,
        TNLean.PEPS.translatedTwoPlaquetteRegion_eq_map,
        TNLean.PEPS.translatedTwoPlaquetteRegion_eq_arcRectangle,
        TNLean.PEPS.translatedTwoPlaquetteRegion_connected,
        TNLean.PEPS.translatedTwoPlaquetteRegion_card_internalEdges,
        TNLean.PEPS.translatedTwoPlaquetteCycleBond_injective,
        TNLean.PEPS.translatedTwoPlaquetteCycleBond_eq_up}
    \leanok
    \uses{def:peps_translated_two_plaquette_geometry,
        thm:peps_two_plaquette_graph,thm:peps_two_plaquette_torus_geometry}
    The actual region $R_v$ is the cyclic rectangle
    $v+\bigl(\{0,1,2\}\times\{0,1\}\bigr)$, with six sites and seven internal bonds.
    Its induced graph is connected and the transported path is a spanning tree.
    The two distinct selected non-tree bonds are the native upward steps
    \begin{align}
        v\longrightarrow v+(0,1),\qquad
        v+(1,0)\longrightarrow v+(1,1).
    \end{align}
    Their endpoint pairs retain the actual ordered-edge convention, which may
    reverse the displayed direction at the vertical seam.
    This is the arbitrary-position finite-torus geometry for
    \cite[Theorem~6.16]{Schuch2010PEPS}; the remaining unrestricted lattice
    scope is recorded in \cite{gap:rmp_peps_quantum_double_g_isometry}.
\end{theorem}

\begin{proof}\leanok
    Torus translation is a graph automorphism. Restricting it to the six-site
    set gives the actual induced graph identification, so it preserves
    the spanning tree, connectedness, and both cardinalities.
    Applying it to the original site coordinates yields the cyclic rectangle.
    The two remaining edges are the translated vertical steps. Sorting their
    endpoint pairs gives exactly the native upward bonds, even at a seam;
    the edge bijection preserves their distinctness.
\end{proof}
